\section{Practical guidance}\label{sec:guidance}

This section gives a procedure for comparing a small set of candidate classifiers on fairness and accuracy with a limited labelling budget.

\paragraph{Step 1: describe the agreement structure before labelling.} Compute $K_i$ and the stratum sizes $N_h$ on the unlabeled pool, and the pairwise disagreement rates of the models to be compared. These numbers cost no labels and indicate the benefit of the design (rank correlation $-0.58$ to $-0.71$ with the efficiency at $L=250$ to $1000$, law school excluded, Section~\ref{sec:empirical}): in our data, disagreement rates below 5\% went with efficiencies of 5.4 to 15.0 at $L=1000$ (law school excluded), while rates of 8.6\% went with efficiencies of 3.8 to 5.0. If the candidate models disagree on a large share of the pool, the gain is modest.

\paragraph{Step 2: check that the minority cells are not empty.} Every group-wise rate needs positive and negative instances in each group. The proxy probabilities give an expected count of positives per group and stratum before any label is collected; if the expected number of minority positives in the pool is very small, no design can estimate the corresponding rates, and the estimand should be reconsidered (for example by merging groups).

\paragraph{Step 3: choose the estimands and allocate.} Use the signed contrasts \eqref{eq:theta} for the pairs of interest. If only a few pairs matter, restricting the estimand set to them concentrates the labels further. Compute the allocation \eqref{eq:alloc} with the average predicted probability as proxy, at least ten labels per stratum, draw the instances in each stratum without replacement and label them.

\paragraph{Step 4: estimate and report uncertainty.} Estimate the totals with the difference estimator \eqref{eq:DIF}, which corrects the pool-wide proxy total by the labelled residuals and needs no tuning, and report the contrasts with standard errors and intervals. If the statement is about the deployment population rather than about the pool itself, add the pool-sampling variance (Section~\ref{sec:population}): the pool-conditional intervals undercover once the labelling budget is a sizeable fraction of the pool, with coverage of the population contrast down to 0.79 when the variance ratio exceeds one, while the two-phase intervals keep it near 0.95. For budgets below about 500 labels or pools below a few thousand instances, use the variance floor of Section~\ref{sec:ablations} (and, if worst-contrast coverage matters, its conservative version, which widens the intervals by about a quarter in the median) and report the stratum sample sizes, so that readers can judge how many minority positives informed each rate.

\paragraph{Step 5: test and decide.} State the hypotheses before labelling: equality of the gaps of two models, or non-inferiority of a candidate against an incumbent with a margin $\delta$ chosen for the application. For a single pre-specified contrast use the Wald test of Section~\ref{sec:testing}. When many contrasts are examined, use the Holm or the max-$t$ procedure, which control the familywise error rate asymptotically; in our simulations max-$t$ had familywise error slightly above nominal (0.058 at the 5\% level) and, on real pools, simultaneous coverage below that of Bonferroni at the smallest budgets (Section~\ref{sec:testing-sim}), so Holm is the safer default when the budget is small. A contrast whose test rejects supports a statement that one model has a smaller gap than the other on the stated metric. When none rejects, the budget has not resolved the comparison. To plan the budget, use the sample-size formula of Section~\ref{sec:testing}: for a fixed allocation the labels needed for a given power are inversely proportional to the squared effect, and the design effect of Section~\ref{sec:empirical} converts the number of simple random labels into the number needed here; if some contrasts matter more, the power-optimal allocation of Section~\ref{sec:testing} favours them.

\paragraph{When not to use the design.} The design compares fixed models. It does not estimate a single model's fairness more precisely than a group-stratified design does, because one model's group-wise rates do not depend on disagreements with other models. It also assumes that the group attribute is observed for the whole pool. When the attribute is missing for some instances, it can be added as a third stratum level, but the fairness estimands then refer only to instances with a recorded group.
